\chapter*{Abstract}
\addcontentsline{toc}{chapter}{Abstract}

From early 2000s, when the term Big Data was first coined, the evolution of technology has imposed new challenges on the design of data systems. The growing amount of data highlighted the limits of traditional architectures, therefore research moved its focus on new solutions able to handle it. Big data are also a remunerative market. It has been measured that, in 2013, the vendors revenue derived by related hardware, software and services was more than 18 million dollars, with a 58\% boost respect to the previous year \cite{bib29}. It is estimated that by 2017 the revenue will be more than 50 million dollars. Recently, we could notice an evolution for cloud-based big data services for large-scale analytics. More and more companies realized about the benefits of data analysis in order to carve out advantages in the market. A survey of 2013 \cite{bib32} shows that more than 50\% of the survey respondents declared that investing on analytics improved organization's competitive position, in terms of identifying new ways to perform sales and understanding users behavior for targeting them with ad-hoc campaigns.\\
Traditionally, for processing a high volume of data, a batch approach with tools like Hadoop is used, specially after the release of Yarn \cite{bib24}, which has broadened its applicability domains. However, some use cases, require to process data in nearly real-time, in order to have the possibility of taking business decisions quickly according to their results. What normally happens, is that these two ways of processing data are independent and are implemented in two different systems.\\
In response to a need of new approaches and patterns for big data processing, this thesis proposes a new architectural pattern: ``VimondAnalytics", named after the company (Vimond Media Solutions \cite{bib1}) where this master thesis project was carried out. It is based on the Lambda architecture, a pattern proposed by N. Marz \cite{bib2}, which provides a new way to combine the classic batch and real-time computations in one solution. Data are analyzed in two layers in parallel, batch and real-time respectively, and their results are stored on dedicated databases and finally merged in the serving layer. High latency of batch results is compensated by the real-time layer. This gives high availability to the system at the cost of eventual consistency according to the CAP theorem \cite{bib33}. Even though Lambda architecture was first introduced in 2011, it has not been deeply investigated up to now, and we could find in literature only few examples of complete implementations \cite{bib30}, \cite{bib31}. The purpose of this thesis is to provide a proof of concept implementation, and show its applicability in the fields of big data analytics. We will show how the initial pattern of Marz can be simplified, by eliminating intermediate databases, performing all the data aggregation and merging directly into the serving layer. The developed system has been deployed on AWS\footnote{https://aws.amazon.com} (Amazon Web Service). \\
The results obtained are straightforward: having a safe data deletion mechanism leads to a decrease in the required time to query the system, compared to a classic solution where all data are ingested into the serving layer and never deleted. In case of throttled input frequency, we will show that VimondAnalytics query time is limited by an upper and a lower bound, depending on the chosen batch time window. We will also show how the system reacts to some failures of single components of the system. However, some of the results, are tightly coupled with the frameworks used for the implementation, and might not be valid outside this context.\\
The rest of this report is structured as follows:
\begin{itemize}
\item Chapter 1 gives an introduction of the context of this work, specifically about big data analytics and what is the state of the art in term of architectures and patterns for handling them. 
\item Chapter 2 briefly analyzes what are the tools and technologies used in this work, highlighting their main features and their architectures.
\item Chapeter 3 describes VimondAnalytics architecture in detail. Great attention is given to each different layer and it is explained how they all cooperate in order to have a working solution.
\item Chapter 4 discusses a possible implementation of VimondAnalytics using the tools introduced in chapter 2.
\item Chapter 5 shows the results obtained with the proposed framework. General purpose tests were run in order to validate these results. For completion, the system was also tested under some common failover scenarios.
\item Finally, chapter 6, summaries the work carried out in this thesis project. Furthermore, it outlines some future work scenarios in terms of enhancement of the current system features and applicability domains. 
\end{itemize}
\vspace{\fill}
\textbf{Keywords:}
Big data, analytics, Lambda architecture, batch processing, real-time processing, cloud computing, distributed systems